\documentclass[11pt]{article}

\setlength{\parindent}{0pt}
\usepackage{xltxtra}
\usepackage{hyperref}
\hypersetup{hidelinks}
\usepackage{url}
\urlstyle{tt}
\usepackage{xcolor}
\definecolor{CVBlue}{RGB}{23,110,191}
\usepackage{calc}
\usepackage{graphicx}
\usepackage{tikz}
\usetikzlibrary{calc}
\usepackage{fontspec}
\usepackage{xeCJK}
\usepackage{enumitem}
\CJKsetecglue{}

\setmainfont[
    Path = fonts/Main/,
    Extension = .otf,
    BoldFont = texgyretermes-bold.otf,
]{texgyretermes-regular.otf}

\setCJKmainfont[
    Path = fonts/hansans/,
    Extension = .ttf,
    BoldFont = NotoSansSC-Bold.ttf,
]{NotoSansSC-Regular.otf}

\usepackage{relsize}
\usepackage{xspace}
\usepackage{fontawesome}

\usepackage[
    a4paper,
    left=1.2cm,
    right=1.2cm,
    top=1.5cm,
    bottom=1cm,
    nohead
]{geometry}

\renewcommand{\baselinestretch}{1.5}

\usepackage{titlesec}
\usepackage{enumitem}
\setlist{noitemsep}
\setlist[itemize]{topsep=0.25em, leftmargin=*}
\setlist[enumerate]{topsep=0.25em, leftmargin=*}

\newlength{\interProjectSpacing}
\setlength{\interProjectSpacing}{0.9em}
\newcommand{\projectsep}{\vspace{\interProjectSpacing}}

\newlength{\intraProjectTitleSep}
\setlength{\intraProjectTitleSep}{0.4em}
\newcommand{\titlebreak}{\\[\intraProjectTitleSep]}

\newlength{\intraProjectListTopSep}
\setlength{\intraProjectListTopSep}{0.2em}

\titleformat{\section}
{\large\bfseries\raggedright}
{}{0em}
{}
[{\color{CVBlue}\titlerule}]
\titlespacing*{\section}{0cm}{*1.6}{*1.2}


\begin{document}
\pagenumbering{gobble}

\begin{tikzpicture}[remember picture, overlay]
    \node[anchor = north east] at ($(current page.north east)+(-2cm,-1.2cm)$) {\includegraphics[height=3cm]{avatar.jpg}};
\end{tikzpicture}%
\begin{tikzpicture}[remember picture, overlay]
    \node[anchor = north west] at ($(current page.north west)+(0.5cm,+1.0cm)$) {\includegraphics[height=6cm]{zju.png}};
\end{tikzpicture}%

\centerline{\LARGE\bfseries{张宇航}}
\centerline{\normalsize{\faPhone\ 130-9832-9363 \quad \faEnvelopeO\ \href{mailto:yuhang3.25@intl.zju.edu.cn}{yuhang3.25@intl.zju.edu.cn}}}


\section{\makebox[\widthof{\faGraduationCap}][c]{\color{CVBlue}\faGraduationCap}\ 教育经历}

\textbf{浙江大学} \quad 国际联合学院 \quad 结构工程 \hfill 2025.09 -- 至今\\[0.3em]
\begin{itemize}[nosep]
    \item 主修课程：具身智能，人工智能算法与系统，人工智能理论和高级应用，AI 赋能：创新思维引领下的产品设计
\end{itemize}

\vspace{0.4em}
\textbf{长安大学} \quad 未来交通学院 \quad 土木工程 \hfill 2021.09 -- 2025.06


\section{\makebox[\widthof{\faUsers}][c]{\color{CVBlue}\faUsers}\ 项目经历}

% --- 项目一 ---
\textbf{全离线 RAG 智能座舱语音交互系统} \hfill 2024.XX -- 2024.XX \titlebreak
项目简介：针对无网络环境（车载/工业场景），设计并实现全离线智能语音交互系统。采用微服务架构将 ASR、RAG、LLM、TTS 四个模块解耦为独立进程，通过 ZeroMQ 异步通信实现语音输入→知识检索→大模型推理→语音播报的完整闭环。
\begin{itemize}[nosep, topsep=\intraProjectListTopSep]
    \item \textbf{流式 ASR}：集成 Sherpa-ONNX + Zipformer 流式模型，基于 PortAudio 异步回调采集音频；调整 VAD 端点检测阈值（默认 1.2s→0.4s）降低首响应时延；设计 ASR-TTS 互斥机制避免 TTS 播报时的回声误触发
    \item \textbf{端侧 RAG}：通过 Pybind11 嵌入 Python 调用 text2vec 实现向量检索；部署规则引擎进行查询分类（紧急/事实/创意），紧急指令直通知识库跳过 LLM 推理实现快速响应
    \item \textbf{LLM 推理}：基于 Ollama 部署 DeepSeek-R1 1.5B（Q4 量化）模型，通过 libcurl 调用 HTTP Streaming API 实现流式推理；设计按中文标点分句的伪流式传输协议，每完成一句立即发送 TTS，降低用户感知的首句等待时间
    \item \textbf{并发 TTS}：基于生产者-消费者模型构建文本/音频双缓冲队列，将语音合成与音频播放分离到独立线程实现流水线并行，消除 LLM 逐句输出时的语音卡顿
    \item \textbf{通信与工程化}：封装 ZeroMQ REQ-REP 为三层继承体系（基类 + Server/Client），四模块独立进程实现故障隔离；全程 RAII + unique\_ptr 管理资源，CMake 统一构建多模块依赖
\end{itemize}

\projectsep

% --- 项目二 ---
\textbf{基于 A2A 协议的 C++ 多 Agent 通信 SDK} \hfill 2025.10 -- 2025.12 \titlebreak
项目简介：基于 C++17 独立实现 Google A2A（Agent-to-Agent）协议的完整 SDK，利用阿里百炼平台提供的 Agent 构建多 Agent 通信系统验证协议准确性。同时利用 Redis 分布式存储对话历史保证上下文记忆，通过 tcpdump 抓包与 Wireshark 分析验证通信流程与协议格式的正确性。
\begin{itemize}[nosep, topsep=\intraProjectListTopSep]
    \item \textbf{A2A 协议实现}：使用 JSON-RPC 2.0 构建 A2A 通信消息格式，支持 message/send（一次性返回）与 message/stream（流式逐步返回）两种消息模式；定义完整数据模型体系，包括信息模型（Agent 身份与能力描述）、消息模型（单轮/多轮对话消息传递）、任务模型（长时任务状态追踪与多次返回）
    \item \textbf{服务注册发现与意图路由}：实现注册中心，Agent 启动时主动注册元信息与能力标签，支持按标签快速查找与轮询负载均衡，后台心跳线程定期健康检查自动剔除故障节点；Orchestrator 接收用户请求后调用 LLM 进行意图识别，自动路由至对应 Agent 处理
    \item \textbf{Client/Server 分层接口}：a2a\_client 提供客户端 API，支持同步/异步消息发送与 SSE 流式轮询；a2a\_server 提供任务生命周期管理与状态追踪存储，定义 ITaskStore 抽象接口并基于 Redis 扩展实现分布式 TaskStore，支持跨进程对话历史共享与上下文记忆
    \item \textbf{多 Agent 编排落地}：基于阿里百炼平台部署 Orchestrator 与 Math Agent 为独立进程，完成"用户请求→意图识别→服务发现→Agent 调用→结果存储→返回用户"完整闭环；通过 contextId 实现跨进程多轮对话上下文连续推理
\end{itemize}


\section{\makebox[\widthof{\faTrophy}][c]{\color{CVBlue}\faTrophy}\ 荣誉奖项}
\begin{itemize}[nosep]
    \item 美国大学生数学建模大赛 \textbf{国际二等奖}
    \item 第七届全国高校智能交通创新与创业大赛 \textbf{国家级一等奖}
    \item 第 15 届中国大学生计算机设计大赛 \textbf{国家级三等奖}
    \item 第 16 届全国大学生数学竞赛 \textbf{省级一等奖}
\end{itemize}


\section{\makebox[\widthof{\faCogs}][c]{\color{CVBlue}\faCogs}\ 专业技能}
\begin{itemize}[nosep]
    \item 熟练掌握 C++ 11，熟悉面向对象三大特性、STL 常用容器，熟悉智能指针、RAII、右值引用与移动语义、Lambda 表达式等 Modern C++ 特性
    \item 熟悉 TCP/IP 协议栈，掌握 TCP 三次握手/四次挥手、TCP 与 UDP 区别、滑动窗口与拥塞控制；了解 IO 多路复用（select/poll/epoll）及 Socket 编程
    \item 熟悉 Linux 系统编程，理解进程与线程区别、虚拟内存机制、进程间通信方式
    \item 了解 RAG（检索增强生成）技术，了解 Embedding 向量化、余弦相似度检索等相关实现
    \item 了解设计模式（生产者-消费者、工厂模式、RAII、单例模式），有实际项目落地经验
\end{itemize}

\end{document}
